\documentclass[11pt]{article}
\usepackage[margin=0.9in]{geometry}
\usepackage[T1]{fontenc}
\usepackage{lmodern}
\usepackage{amsmath,amssymb,amsthm,graphicx,microtype}
\usepackage[hidelinks]{hyperref}

\newtheorem{theorem}{Theorem}
\newtheorem{lemma}[theorem]{Lemma}
\newcommand{\R}{\mathbb R}
\newcommand{\N}{\mathbb N}
\newcommand{\cM}{\mathcal M}
\newcommand{\spn}{\mathfrak{sp}}

\title{Corollary 2.8 Extends to All\\
Same-Generator Power Modules}
\author{Candidate full solution packet\\
\small arXiv:1610.01659, Question 2.10}
\date{11 August 2026}

\begin{document}
\maketitle

\noindent\textbf{Status.} Candidate full affirmative answer, subject to
expert review.  The result follows by applying the source's own representing-
measure theorem and observing that the character positivity set of a power
module depends only on its generators, not on the even exponent.

\section{Source question}

Let $(A,\tau)$ be a locally multiplicatively convex unital commutative real
algebra.  For a fixed family $G=\{g_j:j\in J\}\subset A$, define the
$2d$-power module
\[
 M_{2d}(G):=\sum A^{2d}
 +\left\{\sum_{j\in F}g_j\sigma_j:
 F\subset J\text{ finite},\ \sigma_j\in\sum A^{2d}\right\}.
 \tag{1}
\]
Infusino and Kuhlmann ask in Question 2.10, printed on page 7
of \cite{IK}, whether their Corollary 2.8 extends to the modules
$M_{2d}(G)$ when the same nonempty generator set $G$ is used for every $d$.

\begin{center}
\includegraphics[width=0.94\textwidth]{figures/open_problem_crop.png}
\end{center}

Explicitly, for a $\tau$-continuous linear functional $L:A\to\R$, they ask
whether the following are equivalent:
\begin{enumerate}
\item $L$ is nonnegative on $M_2(G)$;
\item $L$ is nonnegative on $M_{2d}(G)$ for every $d\in\N$;
\item $L$ is nonnegative on $M_{2d}(G)$ for some $d\in\N$.
\end{enumerate}

\section{Proof intuition}

The easy direction is purely algebraic: every $2d$-th power is a square, so
$M_{2d}(G)\subset M_2(G)$.  For the reverse direction, suppose positivity is
known for one exponent.  Proposition 2.5 of the source turns continuity and
$M_{2d}(G)$-positivity into a nonnegative representing measure supported on
the character positivity set of that module.

That support set is independent of $d$.  Indeed, it is simply the set of
characters at which every generator $g_j$ is nonnegative.  Every element of
$M_2(G)$ is pointwise nonnegative there.  Integrating therefore gives
$M_2(G)$-positivity.  This is exactly the missing implication.

\section{The exponent-independent positivity set}

For a power module $M\subset A$, write
\[
 X_M:=\{\alpha\in X(A):\alpha(m)\ge0\text{ for all }m\in M\},
\]
where $X(A)$ is the real character space of $A$.

\begin{lemma}\label{lem:support}
For every $d\in\N$ and every generator family $G$,
\[
 X_{M_{2d}(G)}
 =K_G:=\{\alpha\in X(A):\alpha(g_j)\ge0\text{ for every }j\in J\}.
 \tag{2}
\]
In particular, this set does not depend on $d$.
\end{lemma}

\begin{proof}
If $\alpha\in X_{M_{2d}(G)}$, then
$g_j=g_j1^{2d}\in M_{2d}(G)$, so $\alpha(g_j)\ge0$ for every $j$.
Conversely, suppose $\alpha\in K_G$.  For every $a\in A$,
$\alpha(a^{2d})=\alpha(a)^{2d}\ge0$, and
\[
 \alpha(g_ja^{2d})=\alpha(g_j)\alpha(a)^{2d}\ge0.
\]
Linearity then makes $\alpha$ nonnegative on every finite sum in (1), proving
$\alpha\in X_{M_{2d}(G)}$.
\end{proof}

\section{Main theorem}

\begin{theorem}\label{thm:main}
Let $(A,\tau)$ be an lmc unital commutative real algebra, let
$G\subset A$ be any generator family (possibly infinite), and let
$L:A\to\R$ be $\tau$-continuous.  Then conditions (1)--(3) listed above are
equivalent.  Hence Question 2.10 of \cite{IK} has an affirmative answer; in
fact the proof also covers the empty generator family.
\end{theorem}

\begin{proof}
For every $d$, the inclusion
\[
 \sum A^{2d}\subset\sum A^2
\]
holds because $a^{2d}=(a^d)^2$.  Multiplication by each common generator and
taking finite sums preserves the inclusion, so
$M_{2d}(G)\subset M_2(G)$.  This proves (1)$\Rightarrow$(2), while
(2)$\Rightarrow$(3) is immediate.

Assume (3), and fix $d$ for which $L$ is $M_{2d}(G)$-positive.  Proposition
2.5 of \cite{IK} gives a nonnegative representing measure $\mu$ for $L$ with
\[
 \operatorname{supp}\mu
 \subset X_{M_{2d}(G)}\cap\spn(\tau)
 \subset K_G,
 \tag{3}
\]
where the last equality is Lemma~\ref{lem:support}.  Every
$m\in M_2(G)$ has $\widehat m(\alpha)=\alpha(m)\ge0$ on $K_G$, again by the
lemma with $d=1$.  Therefore
\[
 L(m)=\int_{X(A)}\widehat m(\alpha)\,d\mu(\alpha)\ge0
 \qquad(m\in M_2(G)).
\]
Thus (3)$\Rightarrow$(1), completing the equivalence.
\end{proof}

\section{Verification and novelty}

The proof has only two substantive checks.  First, the module convention in
the source permits arbitrary generator families but only finite sums in each
element; this is precisely (1), and Lemma~\ref{lem:support} applies without
any cardinality assumption.  Second, source Proposition 2.5 applies to any
$2d$-power module of an lmc algebra, not only to the empty-generator case.
The support inclusion it supplies is exactly (3), after which positivity on
$M_2(G)$ is pointwise.

A bounded novelty search through 11 August 2026 covered the run indexes, the
exact wording of Question 2.10, title and author searches, subsequent papers
by the authors on infinite-dimensional moment problems, and combinations of
``same generators,'' ``$2d$-power module,'' and ``Corollary 2.8.''  No later
paper explicitly answering Question 2.10 was found.  The proof is a direct
deduction from results printed immediately before the question; the novelty
claim should therefore be treated cautiously even though no explicit answer
was retrieved.

\medskip
\noindent\textbf{Human-review recommendation.} Review the exact hypotheses of
source Proposition 2.5 and the module definition.  Under the printed
statements, the argument is complete and needs no additional assumption such
as finite generation, Archimedeanness, or completeness.

\begin{thebibliography}{9}
\bibitem{IK}
M. Infusino and S. Kuhlmann,
\emph{Infinite dimensional moment problem: open questions and applications},
Contemp. Math. \textbf{697} (2017), 187--201; arXiv:1610.01659.
\end{thebibliography}

\end{document}
